% Avsnitt 12.2, ur lectures/L12/appendix/b_bit_timer.md.

\section{Att konstruera \code{bit_timer}}
\label{c12:sec:bittimer}

\code{bit_timer} delar 50 MHz-systemklockan i CAN-bitperioder och producerar två encykelspulser som
resten av kontrollern går på: en som markerar var i biten bussen ska samplas, en som markerar bitens
slut.

Varför en bitperiod alls delas upp, och varför sampelpunkten ligger sent i den snarare än mitt i,
står i \canref{5}; det här avsnittet bygger den modul som gör det den kapitlet beskriver, med en
fast sampelpunkt i stället för konfigurerbara segment.

\subsection{Gränssnitt}
\label{c12:sub:btiface}

\modulefig[0.62]{lectures/L12/appendix/images/bit_timer.png}

\begin{booktable}{@{}p{5.6em}p{3.4em}p{5.2em}L@{}}
  \thead{Port} & \thead{Riktn.} & \thead{Typ} & \thead{Betydelse} \\
  \midrule
  \code{clock} & in & \code{std_logic} & 50 MHz systemklocka. \\
  \code{reset_s2_n} & in & \code{std_logic} & Aktiv låg reset, synkroniserad med två vippor. \\
  \code{enable} & in & \code{std_logic} & Kör räknaren medan \code{'1'}; håller den medan
    \code{'0'}. \\
  \code{resync} & in & \code{std_logic} & Tvingar räknaren tillbaka till tick 0 på nästa
    klockflank. \\
  \code{sample} & out & \code{std_logic} & Flagga som anger att bussen ska samplas. \\
  \code{bit_done} & out & \code{std_logic} & Flagga som anger att den här biten är klar och att
    nästa börjar. \\
\end{booktable}

Inga generics. \code{bit_timer_tb} binder sina portar \textbf{positionellt}, så ordningen och typerna
ovan måste stämma exakt. Namnen är dina; att behålla de här får boktexten och din kod att tala om
samma signaler.

Ur \code{can_def} (\chapref{c10:ch}) behövs \code{TICKS_PER_BIT} (= 50, ticken i en bitperiod) och
\code{SAMPLE_TICK} (= 35, alltså 70~\% in, sent nog för att bussen ska ha hunnit stabilisera sig).

\subsection{Beteende}
\label{c12:sub:btbehav}

En enda fritt löpande räknare över \code{0 .. TICKS_PER_BIT - 1} räcker. Vid varje klockflank: ge
båda pulserna förvalet \code{'0'} och sedan, i prioritetsordning:
\begin{enumerate}
  \item \textbf{\code{reset_s2_n = '0'}}: nollställ räknaren och båda pulserna.
  \item \textbf{\code{resync = '1'}}: starta om bitperioden; sätt räknaren tillbaka till \code{0}.
    Det här vinner över \code{enable}, och används en gång per ram, vid SOF, så att den här nodens
    bittajming börjar om från början, i linje med den flank som inleder ramen i stället för var en
    tidigare räkning råkade stå.
  \item \textbf{\code{enable = '1'}}: pulsa \code{sample} under det enda tick där
    \code{counter = SAMPLE_TICK}. Vid \code{counter = TICKS_PER_BIT - 1}, sätt \code{bit_done} och
    låt räknaren slå runt till \code{0}; annars öka den.
  \item \textbf{Annars (avstängd)}: håll räknaren, båda pulserna låga.
\end{enumerate}

Det här är samma idiom med räknare och prioriterade grenar som \chapref{c7:ch}:s \code{timer}, bara
med två pulsutgångar i stället för en.

\begin{aside}
En notering om \code{enable}, eftersom det är den enda porten vars slutliga användning inte är
uppenbar. Den naturliga gissningen är att \code{can_controller} håller timern medan bussen är ledig,
och det visar sig att den inte gör det: \chapref{c16:ch}:s regel för ledig buss räknar recessiva
\textbf{bitperioder}, så timern måste fortsätta gå även i \code{STATE_IDLE}, och den färdiga designen
binder \code{enable} högt i varje tillstånd. Bygg porten ändå; testbänken kontrollerar den, och en
timer som inte går att stoppa är ett sämre byggblock än en som går.
\end{aside}

\subsection{En bitperiod, tick för tick}
\label{c12:sub:bttrace}

Reglerna ovan är lättare att lita på när de setts på en tidsaxel. Med \code{TICKS_PER_BIT = 50} och
\code{SAMPLE_TICK = 35} ser en aktiverad bitperiod ut så här. Varje kolumn är en stigande klockflank;
\code{counter} är värdet processen ser vid den flanken, och de två raderna under den är pulserna som
den flanken producerar.

\begin{textdrawing}
counter      0   1   2  ..  34  35  36  ..  48  49   |   0   1  ..
sample       0   0   0  ..   0   1   0  ..   0   0   |   0   0  ..
bit_done     0   0   0  ..   0   0   0  ..   0   1   |   0   0  ..
\end{textdrawing}

Tre saker att läsa ur det. Båda utgångarna är \textbf{ett tick breda}, inte nivåer. \code{sample}
kommer vid tick 35 och \code{bit_done} vid tick 49, så \textbf{\code{sample} slår alltid till först,
14 tick tidigare i samma period}; en mottagare läser bussen 70~\% in och först därefter tar perioden
slut. Och periodgränsen är omslaget från 49 tillbaka till 0, vilket är samma flank som
\code{bit_done} markerar.

\textbf{Nu samma bild med en \code{resync}.} Anta att timern är 20 tick in i en period när SOF
anländer och \code{can_controller} pulsar \code{resync}:

\begin{textdrawing}
tick        19  20  21  22  ..  55  56  ..  69  70  71
resync       0   1   0   0  ..   0   0  ..   0   0   0
counter     19  20   0   1  ..  34  35  ..  48  49   0
sample       0   0   0   0  ..   0   1  ..   0   0   0
bit_done     0   0   0   0  ..   0   0  ..   0   1   0
\end{textdrawing}

Räknaren kastas mitt i räkningen och startar om, så \code{counter} och \code{tick} stämmer inte
längre överens. Det som spelar roll är följden: \code{bit_done} anländer vid tick \textbf{70}, en hel
\code{TICKS_PER_BIT} efter omsynkroniseringen, snarare än vid tick 49 där den avbrutna perioden
skulle ha slutat. Noden har antagit sändarens bitgräns i stället för sin egen, vilket är hela poängen
med \code{resync} och exakt vad testbänkens fjärde fall mäter.

Lägg också märke till att \code{resync} startar om perioden trots att \code{enable} är hög hela
tiden; det är därför den står över \code{enable} i prioritetsordningen.

\subsection{Vad testbänken låser fast}
\label{c12:sub:bttb}

\begin{itemize}
  \item Hållen i reset stannar \code{sample} och \code{bit_done} på \code{'0'}.
  \item Medan \code{enable = '0'} går räknaren aldrig framåt; båda pulserna stannar på \code{'0'}.
  \item Medan den är aktiverad pulsar \code{sample} exakt vid \code{SAMPLE_TICK} och \code{bit_done}
    exakt vid det sista ticket, en gång per period.
  \item Efter en \code{resync} mitt i en period kommer nästa \code{bit_done} exakt
    \code{TICKS_PER_BIT} tick senare, inte där den avbrutna perioden skulle ha tagit slut.
\end{itemize}

\subsection{Att lägga in den i \code{can_controller}}
\label{c12:sub:btwire}

\code{bit_timer} sällar sig till de två \code{meta_prev}-instanserna från \secref{c12:sec:metaprev}
inuti toppnivån som skrevs i \chapref{c11:ch}, och är det första blocket där inne som vet något om
CAN. Att instansiera den är samma tvådelade ändring i \file{controller/can_controller.vhd} som
\code{meta_prev} var, och den tar \code{reset_s2_n} från synkroniseraren du kopplade in nyss snarare
än porten \code{reset_n}.

Deklarera först de fyra signaler som instansen kopplas till, i arkitekturens deklarativa del:

\begin{vhdlcode}
    -- bit_timer.
    signal bt_enable, bt_resync, bt_sample, bt_bit_done : std_logic;
\end{vhdlcode}

Instansiera sedan själva entiteten, efter \code{begin}:

\begin{vhdlcode}
    bit_timer1: entity work.bit_timer
        port map(clock, reset_s2_n, bt_enable, bt_resync, bt_sample, bt_bit_done);
\end{vhdlcode}

Ingenting driver \code{bt_enable} eller \code{bt_resync} än, och ingenting läser \code{bt_sample}
eller \code{bt_bit_done}; tillståndsmaskinen som gör bådadera är \chapref{c16:ch}:s. Det är väntat i
det här läget. Det som spelar roll nu är att \file{can_controller.vhd} fortfarande analyserar rent,
vilket \code{make build} kontrollerar så snart filen finns, och att \code{bit_timer} själv passerar
\code{bit_timer_tb}.
